\section{Introducción}
    La ciencia de datos es la disciplina que combina estadística, computación y
    conocimiento del dominio para extraer, a partir de datos observados,
    conclusiones que respalden decisiones. Un proyecto de ciencia de datos recorre
    un ciclo que comienza con la formulación de una pregunta, continúa con la
    obtención, descripción y exploración de los datos, sigue con la preparación y
    el modelado, y termina con la evaluación y la comunicación de los resultados
    \cite{chapman2000crisp}, \cite{unisalle2026proyecto}. Este informe corresponde
    a la primera fase de ese ciclo, el planteamiento del problema y el análisis
    exploratorio de datos, cuyo propósito es entender la estructura, la calidad y
    las regularidades del conjunto antes de construir cualquier modelo
    \cite{tukey1977exploratory}.

    \vspace{10pt}

    El caso de estudio es la prueba Saber 11, el examen de Estado que presentan en
    Colombia los estudiantes que terminan la educación media y cuyo puntaje global
    condiciona el acceso a la educación superior, a programas de becas y a los
    cupos de las universidades públicas \cite{icfes2022guia}. El Instituto
    Colombiano para la Evaluación de la Educación (ICFES) publica, en el Portal
    Nacional de Datos Abiertos, un registro anónimo por estudiante que reúne el
    resultado de la prueba y un cuestionario sociodemográfico sobre su hogar y su
    colegio \cite{icfes2023resultados}. La literatura sobre los determinantes del
    rendimiento en Saber 11 ha documentado de forma consistente la asociación del
    puntaje con el nivel socioeconómico del hogar, la educación de los padres, la
    naturaleza y la jornada del colegio y la región \cite{gaviria2001determinantes},
    \cite{chica2010determinantes}, lo que convierte al conjunto en un escenario
    adecuado para un problema de aprendizaje supervisado de regresión en el que la
    pregunta no es solo cuánto se puede predecir, sino qué variables, ajenas al
    esfuerzo individual del estudiante, explican esa predicción.

    \vspace{10pt}

    El trabajo se desarrolla en un cuaderno de Google Colab que descarga el
    periodo 2022-4 del conjunto, lo describe, verifica su calidad, caracteriza la
    variable objetivo y examina su relación con las variables del estudiante, del
    hogar, del colegio y de la geografía. Una de las verificaciones es
    metodológicamente central, puesto que el puntaje global se construye como un
    promedio ponderado de los puntajes por área, estos últimos no pueden usarse
    como variables explicativas sin incurrir en una fuga de información
    \cite{kaufman2012leakage}, y el análisis lo comprueba sobre los datos. El
    informe está organizado de la siguiente manera, la Sección~II plantea el
    problema, la pregunta de investigación, los objetivos y las hipótesis; la
    Sección~III describe el conjunto de datos; la Sección~IV expone la metodología
    del análisis exploratorio; la Sección~V presenta los resultados y los primeros
    hallazgos, y la Sección~VI recoge las conclusiones y las decisiones que
    orientan la fase de preparación y modelado.
